% Generated draft: requires mathematical and source review.
\documentclass[11pt]{article}
\usepackage[margin=25mm]{geometry}
\usepackage{fontspec}
\setmainfont{Latin Modern Roman}
\usepackage{amsmath,amssymb,mathtools,unicode-math}
\setmathfont{Latin Modern Math}
\usepackage{enumitem,needspace,xurl,hyperref}
\setlength{\parindent}{0pt}
\setlength{\parskip}{5pt}
\setlength{\emergencystretch}{3em}
\newcommand{\N}{\mathbb{N}}
\newcommand{\Z}{\mathbb{Z}}
\newcommand{\Q}{\mathbb{Q}}
\newcommand{\R}{\mathbb{R}}
\newcommand{\C}{\mathbb{C}}
\newcommand{\F}{\mathbb{F}}
\newcommand{\A}{\mathbb{A}}
\newcommand{\PP}{\mathbb{P}}
\newcommand{\E}{\mathbb{E}}
\newcommand{\eps}{\varepsilon}
\newcommand{\dd}{\,\mathrm{d}}
\begin{document}
\section*{UnsolvedMath ID: 30001952}
\textbf{Source problem code:} catalogue entry
\subsection{Definitions and mathematical statement}
The supplied existence question is:
\[
\text{Does there exist an algebraic torus quotient }Y=A/G
\text{ by a freely acting finite group }G
\text{ such that }b_2(Y)=1?
\]
An action of a finite group \(G\) on \(A\) is free if
\[
\forall g\in G\;\forall a\in A,\qquad g\cdot a=a\Longrightarrow g=e_G,
\]
where \(e_G\) is the identity element of \(G\). The notation \(A/G\) denotes the quotient by this action. If a complex variety with its usual topology is intended, the second Betti number is
\[
b_2(Y):=\dim_{\Q}H^2(Y;\Q),
\]
where \(H^2\) denotes singular cohomology.

The formulation has unresolved domain qualifications. The supplied record does not define ``algebraic torus'': a complex algebraic group \((\C^\times)^n\) and a projective complex torus (an abelian variety) are different possible readings. The ground field, permitted automorphisms, and intended quotient category are not stated. Consequently, the complex-topological definition of \(b_2\) above is conditional on that interpretation. No dimension restriction appears in the question. The supplied background describes the question as open beyond dimension three, but does not specify a precise dimension quantifier or provide the lower-dimensional conclusions. It therefore does not justify silently replacing the displayed question by an existence question in dimension greater than three. These qualifications require source clarification before a fully specified mathematical statement can be accepted.
\subsection{Short English statement}
The source asks whether a free finite quotient of an object called an algebraic torus can have second Betti number one. The meaning of algebraic torus and the intended dimension range remain unresolved. The supplied curation labels the problem open beyond dimension three; that dated assessment is not independently verified here.
\subsection{Sources}
\begin{itemize}
\item Supplied local UnsolvedMath registry record, permanent ID \(30001952\), title ``Torus Quotients with Second Betti Number One,'' source problem code \texttt{catalogue entry}, in the Oberwolfach Reports Open Problems collection. The record gives proposed year 2012 and names no individual proposer.
\item Supplied original reference: \emph{Mini-Workshop: Manifolds with Lower Curvature Bounds} (2012), Oberwolfach Reports, pp.\ 5--42, DOI \texttt{10.4171/owr/2012/01}. \url{https://doi.org/10.4171/owr/2012/01}.
\item Supplied literature source labeled ``finite three dimensional torus quotients'': \url{https://doi.org/10.5802/aif.3326}. No authors or full bibliographic title are supplied in the record.
\item Literature assessment dated exactly \texttt{2026-08-21}: the supplied curation labels the problem open and states that the original report left this existence question open beyond dimension three, while later work on three dimensional torus quotients does not settle the higher dimensional case. This is supplied dated curation, not an independently verified literature assessment; the record itself requests independent verification of time-sensitive status and cited claims.
\item The registry labels the statement \texttt{corrected\_verified} and reports a manual semantic, typographic, and LaTeX review asserting that the cleanup is self-contained and meaning-preserving. This is supplied review metadata, not independent mathematical verification of the formulation above.
\end{itemize}
\end{document}
